% ============================================================
% Chapter III — Methodology
% ============================================================
\section{Methodology}

\subsection{Research Approach}

This study employs a quantitative, cross-sectional research design using
multivariate statistical methods. The approach is predictive and
confirmatory: features are engineered from within-race telemetry data, the
statistical model is fitted using maximum likelihood estimation (MLE), and
generalizability is assessed through leave-one-group-out cross-validation
(LOGO-CV). The research follows a sequential analytical pipeline---data
ingestion, feature engineering, assumption testing, inferential modeling,
and machine learning validation.

\subsection{Design of the Study}

The study uses a \textbf{multinomial logistic regression} framework with a
three-class ordinal response variable (Finishing Tier: Podium, Point Scorer,
No Points). The design involves:

\begin{enumerate}
  \item \textbf{Assumption testing} via Box's M test to determine the
    appropriate multivariate method.
  \item \textbf{Multicollinearity diagnostics} via VIF to validate feature
    independence.
  \item \textbf{Inferential modeling} via MLE to estimate log-odds
    coefficients.
  \item \textbf{Predictive validation} via LOGO-CV to assess generalizability.
\end{enumerate}

\subsection{Sampling Strategy}

The unit of analysis is the \textbf{individual lap}. Every lap from every
driver across three selected races was included, yielding 2,880 observations
after data cleaning. No random sampling was performed; the dataset represents
the complete population of completed laps from the three selected races.
Races were chosen purposively to cover three architecturally distinct circuit
types:

\begin{enumerate}
  \item \textbf{Bahrain Grand Prix} (Sakhir) --- permanent desert circuit
    with high tire degradation.
  \item \textbf{Saudi Arabian Grand Prix} (Jeddah) --- high-speed street
    circuit with minimal run-off.
  \item \textbf{Australian Grand Prix} (Albert Park) --- parkland circuit
    with mixed corner profiles.
\end{enumerate}

\subsection{Tools and Instruments}

The following tools and libraries were used:

\begin{table}[htbp]
  \centering
  \caption{Software and Libraries}
  \label{tab:tools}
  \begin{tabular}{lll}
    \toprule
    \textbf{Tool} & \textbf{Version} & \textbf{Purpose} \\
    \midrule
    Python & 3.11 & Data preparation and ML validation \\
    FastF1 & 3.8.2 & F1 telemetry data extraction \\
    pandas & 2.3.3 & Data manipulation \\
    scikit-learn & 1.8.0 & Logistic regression, LOGO-CV \\
    R & 4.x & Statistical inference (Box's M, MLE) \\
    nnet (R) & --- & Multinomial logistic regression \\
    car (R) & --- & VIF diagnostics \\
    heplots (R) & --- & Box's M test \\
    \bottomrule
  \end{tabular}
\end{table}

\subsection{Data Collection Process}

Data collection proceeded in five stages:

\begin{enumerate}[label=Phase \arabic*:, leftmargin=*]
  \item \textbf{Session Loading.} Each race session was loaded via
    \texttt{fastf1.get\_session()} with telemetry, lap, and weather data
    enabled. The FastF1 cache system stored downloaded data locally.

  \item \textbf{Multi-Race Stacking.} All laps from the three sessions were
    concatenated into a single dataframe. Each observation was tagged with a
    \texttt{Race\_ID} (e.g., \texttt{2023\_Bahrain}) for group-aware
    cross-validation.

  \item \textbf{Target Variable Definition.} The \texttt{Finishing\_Tier}
    variable was mapped from \texttt{Position}:
    \begin{itemize}
      \item \textbf{Podium:} Position $\leq 3$
      \item \textbf{Point Scorer:} $4 \leq$ Position $\leq 10$
      \item \textbf{No Points:} Position $> 10$
    \end{itemize}

  \item \textbf{DNF Protocol.} Drivers with \texttt{Status = "Retired"} or
    \texttt{ClassifiedPosition = "R"} were flagged, and all their completed
    laps were reassigned to the ``No Points'' tier. This prevents a driver
    running P2 pace before a mechanical failure from contaminating the Podium
    cluster.

  \item \textbf{Feature Engineering.} Three metric predictors were constructed:
    \begin{itemize}
      \item \textit{Fuel Mass Proxy:} Linear burn-off from 110\,kg to 0\,kg:
        \[
          \text{Fuel Mass Proxy} = 110 - \frac{110}{\text{Max Lap}} \times
          \text{Lap Number}
        \]
      \item \textit{Tire Age:} Directly sourced from \texttt{TyreLife}.
      \item \textit{Degradation $\times$ Fuel:} Mean-centered interaction:
        \[
          \text{Deg} \times \text{Fuel} =
          (\text{Tire Age} - \bar{x}_{\text{Tire Age}}) \times
          (\text{Fuel Mass} - \bar{x}_{\text{Fuel Mass}})
        \]
      \item \textit{Tire Compound Dummies:} $k{-}1$ dummy coding with the
        first alphabetical compound (Hard) as baseline.
    \end{itemize}
\end{enumerate}

\subsection{Ethics and Compliance}

This study uses publicly available telemetry data from the FastF1 API, which
interfaces with the official Formula One data feed. No human participants were
involved, and no personally identifiable information was collected. Driver
identifiers (e.g., ``VER'' for Max Verstappen) are public knowledge and used
solely for data grouping. The study complies with FastF1's terms of use for
academic research.

\subsection{Ensuring Credibility}

Several measures were taken to ensure the credibility and reproducibility of
the analysis:

\begin{itemize}
  \item \textbf{Data transparency:} The full dataset (\texttt{f1\_clean\_data.csv})
    and all source code (Jupyter notebooks, R scripts) are available in the
    project repository for independent verification.
  \item \textbf{Assumption testing:} Box's M test and VIF diagnostics were
    conducted before model fitting, and results are reported regardless of
    outcome.
  \item \textbf{Strict validation:} LOGO-CV was chosen over random train-test
    splits to prevent data leakage and produce honest out-of-sample performance
    estimates.
  \item \textbf{Balanced evaluation:} Class weighting was applied during
    validation, and performance is reported via precision, recall, and F1-score
   ---not just accuracy---to prevent misleading claims on imbalanced data.
  \item \textbf{Reproducibility:} The complete analytical pipeline---from raw
    data ingestion through statistical testing to ML validation---is documented
    in executable notebooks with fixed random seeds where applicable.
\end{itemize}
